% Body only: the heading is supplied by toptesi's \sommario in thesis.tex.

Mobile applications used in the field -- on construction sites, in
agriculture, during humanitarian response, or, as in the scenario studied
in this thesis, for facility maintenance reporting on a university campus
-- routinely operate in environments where network connectivity is
intermittent, slow, or entirely absent. Applications that assume an
always-on connection degrade badly in these conditions: user input is
lost, screens hang waiting for a server response, and users lose trust in
the tool. \emph{Offline-first} design addresses this by treating the local
device as the primary source of truth and the network as an optional,
asynchronous channel for synchronization.

This thesis presents the design, implementation, and evaluation of
\emph{CampusFix}, an offline-first mobile application for reporting
maintenance problems (broken projectors, damaged furniture, Wi-Fi outages,
HVAC issues, missing lab equipment) on a university campus. The application
is built with Kotlin Multiplatform and Compose Multiplatform, sharing a
single codebase across Android and iOS, and persists data locally in an
embedded SQLite database accessed through the Room persistence library. A
dedicated synchronization module observes device connectivity, queues
locally-created reports, and pushes them to a Firebase Firestore backend
when a connection becomes available, exposing an explicit per-report sync
status (\texttt{pending}, \texttt{synced}, \texttt{failed},
\texttt{deleted\_pending}) so that users always know whether their data has
left the device.

The thesis makes three contributions. First, it distills the engineering
questions that recur when building offline-first mobile software -- what to
store locally, when to trigger synchronization, how to surface sync state
to the user, and how to resolve conflicting edits -- into a concrete
requirements analysis grounded in a realistic scenario. Second, it presents
a layered system architecture (presentation, domain, data, and
synchronization layers) and a synchronization algorithm based on a
local-first write path with a background push-and-reconcile strategy,
together with a last-write-wins conflict resolution policy suitable for
low-contention, single-owner records such as maintenance reports. Third, it
evaluates the resulting system against six scenarios -- offline creation,
delayed synchronization, failed synchronization, large sync queues,
concurrent edit conflicts, and end-user usability -- discussing both the
quantitative behaviour of the synchronization module and the qualitative
experience of first-time users.

The results indicate that a relatively small, well-isolated synchronization
module is sufficient to make a CRUD-style field-reporting application fully
usable offline, that connectivity-triggered background sync keeps the
perceived latency of ``going back online'' low, and that a strict
last-write-wins policy is an acceptable trade-off for this class of
application, provided the user interface makes the sync status visible at
all times. The thesis closes with a discussion of the current
implementation's limitations -- most notably that synchronization is
presently a one-directional push from device to server rather than a full
bidirectional merge -- and outlines the engineering work required to close
that gap.
